% Pursuing Stacks (cotes 134-2 to 134-8): archive cote 134-7, site page 56 (PDF page 57), typescript page 498.
% typescript. Draft compared: ps6.tex:2771-2829.
% AI-transcribed by Claude Opus 5.5 (claude-opus-5-5, adaptive thinking, effort high), Message Batches API, batch ps4 (msgbatch_012Aoc1NuskATPsyBeKSYK7J), 2026-11-12; not yet checked by eye.
\scanpage{134-7/56}
\mspage{498}

the homology invariants a lot more precise\inserted{presumably,} than those currently used so
far.  Of course,  the significance of (9) for the study of the usual,
``discrete'' homotopy types,  will be subordinated to \unsure{haw|how} difficult it will
turn out for such a homotopy type to be  a) realizable by a schematic
one over $\underline{Z}$,  and  b) to get hold of a\xout{n explicit} more or less explicit
description of \xout{a} such a schematic homotopy type,  via say a semisimplicial
bundle as a model.  One is of course thinking more specifically of the
case of the spheres $S^n$ ,  the first case (besides the trivial $S^1$ case)
being $S^2$ ,  the sequence of homotopy groups of which (as far as I know)
is not understood as yet.  Viewing the spheres $S^n$ as successive suspensi-%
ons of $S^1$ ,  which $S^1$ is fitting nicely into the formalism of schematic
homotopy types as a $K(\underline{Z},1)$  (except that the \inserted{1-connectedness} condition $X_1=e$ \xout{will} \inserted[typed]{is} not \xout{be}
satisfied),  this brings near the question of defining \inserted{the suspension} \deleted{such} operation\deleted{s}
in the \inserted{relevant} derived category $\hand{\mathcal{M}}_o(k)$ or $\hand{\mathcal{M}}_1(k)$\unsure{,} \hand{(}whereas before we had met with
the ``dual'' question of constructing homotopy fibers of maps,  such as
loop spaces\hand{)}.  Thus it would seem that \uline{a breakthrough in getting hold of}
\uline{the standard homotopy constructions within the schematic context},  assu-%
ming that these constructions do still make a sense,  \uline{may well mean a}
\uline{significant advance} \xout{in} \uline{for the understanding of the homotopy groups of}
\uline{spheres}.  This \deleted{at the same time} looks like a very strong motivation for
trying to \xout{develop the} carry through those constructions (possibly even
construct a corresponding ``derivator'' embodying any kinds of finite
``integration'' and ``cointegration'' operations on schematic homotopy types)\hand{,}
and \inserted{(at the same time)} warns us that such work,  if at all feasible,  will most probably be
a highly non trivial one.

The question of``schematizing'' the homotopy types of $S^2$ and $S^3$ reminds
me of a fact which struck me a long time ago (maybe J.P. Serre or someone
else pointed it out to me first).  Namely,  in some respects there does'nt
seem to be \xout{a} really satisfactory algebraic models for these homotopy types\hand{,}
taking into account the basic relationship between the two,  namely:
\over{t}{\unsure{T}}he 2-sphere (or,  equivalently,  the projective complex line)  is a
\uline{homogeneous space} under the quaternionic \uline{group} $S^3$ (or,  equivalently,
under the complex linear group $S\unsure{l}(2,\underline{\unsure{C}})$).  This \xout{sit} relationship,  and
it's manifold ``avatars'' in the realms of discrete groups, Lie groups,
algebraic  groups or group schemes etc,  is one of the \inserted{few} key situations
met with\hand{,} and of equal basic significance\hand{,} in the most diverse quarters\inserted{in mathematics,}
from topology to arithmetics.   Thus $S\unsure{l}(2)$ \xout{plays the r\^ole of \uline{the}} as a
simple Lie (or algebraic) group of minimum rank 1,  plays the  role of

\margin{56}
